% ==============================
% Chapter 1 – Introduction
% ==============================

\chapter{Introduction}

% --------------------------------
% 1.1 Introduction
% --------------------------------

\section{Introduction}

\textbf{Writing Guidelines (400–600 words)}

\begin{itemize}
    \item Provide background of the selected problem domain.
    \item Describe current industry practices and existing solutions.
    \item Identify key challenges, inefficiencies, or gaps.
    \item Justify the need for the proposed system.
    \item Briefly introduce your proposed system and its purpose.
    \item Maintain a formal academic tone and avoid informal language.
\end{itemize}

\vspace{0.5cm}

\textbf{\textit{Replace this guideline text with your final written content before submission.}}



% --------------------------------
% 1.2 Vision Statement
% --------------------------------

\section{Vision Statement}

\textbf{Writing Guidelines (150–250 words)}

\begin{itemize}
    \item Clearly define the target users or stakeholders.
    \item State the core problem the system intends to solve.
    \item Describe the long-term vision and expected impact.
    \item Highlight innovation and uniqueness.
    \item Focus on strategic impact rather than implementation details.
\end{itemize}

\vspace{0.5cm}

\textbf{\textit{Insert your Vision Statement here.}}



% --------------------------------
% 1.3 Related System Analysis
% --------------------------------

\section{Related System Analysis}

\textbf{Writing Guidelines (500–700 words)}

\begin{itemize}
    \item Discuss 2–4 existing systems relevant to your project.
    \item Describe their main features and functionalities.
    \item Critically evaluate their weaknesses or limitations.
    \item Compare them with your proposed solution.
    \item Clearly justify how your system improves upon existing approaches.
\end{itemize}

\vspace{0.5cm}

\textbf{\textit{Insert detailed analytical discussion here.}}

\vspace{0.5cm}

\begin{table}[h!]
\centering
\caption{Comparative Analysis of Existing Systems}
\renewcommand{\arraystretch}{1.3}

\begin{tabularx}{\textwidth}{|p{3cm}|X|X|}
\hline
\textbf{Application Name} & \textbf{Identified Weaknesses} & \textbf{Proposed Improvements} \\
\hline
System A &
\begin{itemize}[leftmargin=*, noitemsep, topsep=0pt]
    \item Weakness 1
    \item Weakness 2
\end{itemize}
&
\begin{itemize}[leftmargin=*, noitemsep, topsep=0pt]
    \item Improvement 1
    \item Improvement 2
\end{itemize}
\\
\hline
System B &
\begin{itemize}[leftmargin=*, noitemsep, topsep=0pt]
    \item Weakness 1
    \item Weakness 2
\end{itemize}
&
\begin{itemize}[leftmargin=*, noitemsep, topsep=0pt]
    \item Improvement 1
    \item Improvement 2
\end{itemize}
\\
\hline

\end{tabularx}
\end{table}



% --------------------------------
% 1.4 Project Deliverables
% --------------------------------

\section{Project Deliverables}

\textbf{Writing Guidelines (400–600 words)}

\begin{itemize}
    \item Clearly list all tangible outputs of the project.
    \item Explain the purpose and functionality of each deliverable.
    \item Identify system modules developed.
    \item Include documentation artifacts (SRS, SDD, test cases, etc.).
    \item Mention deployment or implementation outputs if applicable.
\end{itemize}

\vspace{0.5cm}

\textbf{List of Deliverables:}

\begin{enumerate}
    \item Deliverable 1 – Description of functionality and purpose.
    \item Deliverable 2 – Description of functionality and purpose.
    \item Deliverable 3 – Description of functionality and purpose.
\end{enumerate}



% --------------------------------
% 1.5 System Limitations / Constraints
% --------------------------------

\section{System Limitations / Constraints}

\textbf{Writing Guidelines (200–300 words)}

\begin{itemize}
    \item Identify technical limitations of the system.
    \item Mention hardware or software dependencies.
    \item Discuss time, budget, or resource constraints.
    \item Include scalability or deployment limitations if applicable.
\end{itemize}

\vspace{0.5cm}

\textbf{Example Format:}

\begin{itemize}
    \item Limitation 1 – Brief explanation.
    \item Limitation 2 – Brief explanation.
    \item Limitation 3 – Brief explanation.
\end{itemize}



% --------------------------------
% 1.6 Tools and Technologies
% --------------------------------

\section{Tools and Technologies}

\textbf{Writing Guidelines (200–300 words)}

\begin{itemize}
    \item Justify the selection of each major technology.
    \item Mention backend, frontend, and database tools.
    \item Describe the development environment.
    \item Provide brief technical justification for choices.
\end{itemize}
\begin{table}[h!]
\centering
\caption{Tools and Technologies Used in the Project}
\renewcommand{\arraystretch}{1.3}
\begin{tabularx}{\textwidth}{|p{4cm}|p{3cm}|X|}
\hline
\textbf{Tool / Technology} & \textbf{Version} & \textbf{Purpose} \\
\hline
Tool 1 & Version & Purpose description \\
Tool 2 & Version & Purpose description \\
Tool 3 & Version & Purpose description \\
\hline
\end{tabularx}
\end{table}




% --------------------------------
% 1.7 Relevance to Course Modules
% --------------------------------

\section{Relevance to Course Modules}

\textbf{Writing Guidelines (400–600 words total)}

\begin{itemize}
    \item Demonstrate integration of academic knowledge into practical implementation.
    \item Reflect on how theoretical concepts were applied.
    \item Connect specific course modules to project components.
    \item Highlight achievement of learning outcomes.
\end{itemize}

\vspace{0.5cm}

\subsection{Programming Fundamentals}

\textit{Explain (80–120 words) how programming logic, object-oriented principles, modular design, and code structuring were applied in the project.}


\subsection{Data Structures and Algorithms}

\textit{Explain (80–120 words) how data handling, searching, sorting, optimization techniques, or algorithms were utilized.}


\subsection{Database Management Systems}

\textit{Explain (80–120 words) how database design, normalization, relationships, and queries were implemented.}


\subsection{Software Engineering}

\textit{Explain (80–120 words) how SDLC models, UML diagrams, testing strategies, documentation, and requirement engineering were applied.}


\subsection{Other Relevant Courses}

\textit{Explain (80–120 words) how additional subjects (e.g., AI, Networking, Web Development, Cybersecurity, HCI, etc.) contributed to the project implementation.}

%%%%%%%%%%%%%%%%%%%%%%%%%%%%%%%%%%%%%%%%%%%%%%%%%%%%%%
\section{Chapter Summary}
%%%%%%%%%%%%%%%%%%%%%%%%%%%%%%%%%%%%%%%%%%%%%%%%%%%%%%

This chapter introduced the background and context of the proposed system, 
highlighting existing challenges and the need for improvement within the 
selected problem domain. The vision and strategic objectives of the system 
were outlined to establish its intended impact and innovation. A comparative 
analysis of related systems identified key limitations in current solutions. 
Project deliverables, system constraints, selected technologies, and academic 
relevance were also discussed. 

This chapter establishes the foundation for the detailed problem definition 
presented in the next chapter.
